\subsection{CHAMBERS AND ALCOVES}

Denote by \(t_{r}\) the set of elements \(x \in t\) such that exp x is regular, in other words belongs to \(T_{r}\) . An element x of t belongs to \(t - t_{r}\) if and only if there exists a root \(\alpha \in \mathrm{R}(\mathrm{G}, \mathrm{T})\) such that \(\delta(\alpha)(x) \in 2\pi i\mathbf{Z}\) . For each root \(\alpha \in \mathrm{R}(\mathrm{G}, \mathrm{T})\) and each integer n, denote by \(H_{\alpha,n}\) the set of \(x \in t\) such that \(\delta(\alpha)(x) = 2\pi in\) . The \(H_{\alpha,n}\) are called the singular hyperplanes of t, and \(t - t_{r}\) is the union of the singular hyperplanes. The alcoves of t are the connected components of \(t_{r}\) , and the chambers are the connected components of the complement in t of the union of those singular hyperplanes that pass through the origin (that is, the \(H_{\alpha,0} = \operatorname{Ker}\delta(\alpha), \alpha \in \mathrm{R}(\mathrm{G}, \mathrm{T})\) ).

We have \(\varGamma(\mathrm{T})\subset\mathfrak{t}-\mathfrak{t}_{r}\) ; denote by N(G,T) the subgroup of \(\varGamma(\mathrm{T})\) generated by the nodal vectors (§4, no. 5); by Prop. 11 of §4, no. 6, the quotient \(\varGamma(\mathrm{T})/\mathrm{N}(\mathrm{G},\mathrm{T})\) can be identified with the fundamental group of G.

Finally, denote by W the Weyl group of G relative to T, considered as a group of automorphisms of T and of t, and denote by \(W_{a}\) (resp. \(W_{a}^{\prime}\) ) the group of automorphisms of the affine space t generated by W and the translations \(t_{\gamma}: x \mapsto x + \gamma\) for \(\gamma \in \mathrm{N}(\mathrm{G}, \mathrm{T})\) (resp. for \(\gamma \in \Gamma(\mathrm{T})\) ).

Let \(w \in W\) , \(\gamma \in \Gamma(T)\) , \(\alpha \in R(G, T)\) and \(n \in Z\) . We have:

\[
w (\mathrm{H} _ {\alpha , n}) = \mathrm{H} _ {w \alpha , n}, t _ {\gamma} (\mathrm{H} _ {\alpha , n}) = \mathrm{H} _ {\alpha , n + \langle \gamma , \alpha \rangle}.
\]

It follows that, for all chambers C and all \(w \in W\) , \(w(\mathrm{C})\) is a chamber and that for all alcoves A and all \(w \in W_{a}^{\prime}\) , \(w(\mathrm{A})\) is an alcove. Note that, when \(\mathrm{X}(\mathrm{T}) \otimes \mathbf{R}\) is identified with \(t^{*}\) via the isomorphism \((2\pi i)^{-1}\delta\) , the alcoves of t and the group \(W_{a}\) are the alcoves and the affine Weyl group associated to the root system \(\mathrm{R}(\mathrm{G}, \mathrm{T})\) (Chap. VI, §2, no. 1).

PROPOSITION 2. a) The group \(W_{a}\) (resp. \(W_{a}^{\prime}\) ) is the semi-direct product of W by N(G, T) (resp. \(\Gamma(T)\) ); the subgroup \(W_{a}\) of \(W_{a}^{\prime}\) is normal.

\begin{enumerate}
\def\labelenumi{\alph{enumi})}
\setcounter{enumi}{1}
\item
  The group W (resp. \(\mathrm{W}_a\) ) operates simply-transitively on the set of chambers (resp. alcoves).
\item
  Let C be a chamber and A an alcove. Then \(\overline{C}\) (resp. \(\overline{A}\) , resp. A) is a fundamental domain for the operation of W on t (resp. of \(W_{a}\) on t, resp. of \(W_{a}\) on \(t-t_{r}\) ). If \(x \in t_{r}\) and \(w \in W_{a}\) are such that \(w(x) = x\) , then w = Id.
\item
  For every chamber C, there exists a unique alcove A such that \(A \subset C\) and \(0 \in \overline{A}\) . For every alcove A, there exists a unique \(\gamma \in \mathrm{N}(\mathrm{G}, \mathrm{T})\) such that \(\gamma \in \overline{A}\) .
\end{enumerate}

If \(w \in W\) and \(\gamma \in \Gamma(T)\) , then \(wt_{\gamma}w^{-1} = t_{w(\gamma)}\) and \(wt_{\gamma}w^{-1}t_{\gamma}^{-1} = t_{w(\gamma)-\gamma}\) , with \(w(\gamma) - \gamma \in \mathrm{N}(\mathrm{G}, \mathrm{T})\) ; this immediately implies a). The rest of the proposition follows from Chap. VI, §1, no. 5 and §2, nos. 1 and 2.

COROLLARY 1. Let A be an alcove of t, \(\overline{A}\) its closure, and \(H_{A}\) the stabilizer of A in \(W_{a}^{\prime}\) .

\begin{enumerate}
\def\labelenumi{\alph{enumi})}
\item
  The group \(W_{a}^{\prime}\) is the semi-direct product of \(H_{A}\) by \(W_{a}\) .
\item
  The exponential map \(\overline{\mathrm{A}}\to \mathrm{T}\) and the canonical injection \(\mathrm{T}\rightarrow \mathrm{G}\) induce by passage to the quotients and to subsets homeomorphisms
\end{enumerate}

\[
\begin{array}{l} \overline {{\mathrm{A}}} / \mathrm {H_ {A}} \to \mathrm{T} / \mathrm{W} \to \mathrm{G} / \mathrm{Int} (\mathrm{G}) \\ \mathrm{A} / \mathrm {H_ {A}} \to \mathrm {T_ {r}} / \mathrm{W} \to \mathrm {G_ {r}} / \mathrm{Int} (\mathrm{G}). \end{array}
\]

Let \(w' \in W_a'\) ; then \(w'(A)\) is an alcove of \(t\) , and there exists (Prop. 2 b)) a unique element \(w\) of \(W_a\) such that \(w(A) = w'(A)\) , that is \(w^{-1}w' \in H_A\) . Since \(W_a\) is normal in \(W_a'\) , this proves \(a\) ).

The canonical injection of \(\overline{A}\) into t induces a continuous bijection \(\theta: \overline{A} \to t/W_{a}\) (Prop. 2 c)), and this is a homeomorphism since \(\overline{A}\) is compact. Since \(W_{a}\) is normal in \(W_{a}^{\prime}\) , the group \(H_{A}\) operates canonically on \(t/W_{a}\) (Algebra, Chap. I, §5, no. 4) and \(t/W_{a}^{\prime}\) can be identified with the quotient \((t/W_{a})/H_{A}\) ; the map \(\theta\) is compatible with the operations of \(H_{A}\) , hence induces by passage to the quotient a homeomorphism \(\overline{A}/H_{A} \to t/W_{a}^{\prime}\) . Further, \(\exp_{\Gamma}\) induces a homeomorphism from \(t/\Gamma(T)\) to T, hence also a homeomorphism from \(t/W_{a}^{\prime}\) to T/W. Assertion b) follows from that and Cor. 1 of Prop. 5 of §2, no. 4.

Remarks. 1) The group \(H_{A}\) can be identified naturally with \(\Gamma(\mathrm{T})/\mathrm{N}(\mathrm{G},\mathrm{T})\) , hence also with \(\pi_{1}(\mathrm{G})\) . Thus, it reduces to the identity element when G is simply-connected.

\begin{enumerate}
\def\labelenumi{\arabic{enumi})}
\setcounter{enumi}{1}
\item
  Let \(x \in \mathrm{A}\) ; then \(\exp x \in \mathrm{T}_r\) , so \(\mathrm{Z}(\exp x)_0 = \mathrm{T}\) . Further, \(\exp x\) is very regular (no. 1, Remark) if and only if \(w(x) \neq x\) for all \(w \in \mathrm{W}_a'\) distinct from the identity. By Cor. 1, this also means that \(h(x) \neq x\) for all \(h \in \mathrm{H}_{\mathrm{A}}\) distinct from the identity. In particular, if \(\mathrm{G}\) is simply-connected, then \(\mathrm{Z}_{\mathrm{G}}(t) = \mathrm{T}\) for all \(t \in \mathrm{T}_r\) , and every regular element of \(\mathrm{G}\) is very regular.
\item
  The special points of \(W_{a}\) (Chap. VI, §2, no. 2) are the elements x of t such that \(\delta(\alpha)(x) \in 2\pi i\mathbf{Z}\) for all \(\alpha \in \mathrm{R}(\mathrm{G}, \mathrm{T})\) (loc. cit., Prop. 3), that is such that \(\exp x \in \mathrm{C}(\mathrm{G})\) (§4, no. 4, Prop. 8). For such an element x we have \(wx - x \in \mathrm{N}(\mathrm{G}, \mathrm{T})\) for all \(w \in W\) (Chap. VI, §1, no. 9, Prop. 27), so the stabilizers of x in \(W_{a}\) and \(W_{a}^{\prime}\) coincide. Let S be the set of special points of \(\overline{A}\) ; it follows from the preceding and Cor. 1 that the group \(H_{A}\) operates freely on S, and that the exponential map induces a bijection from \(S/H_{A}\) to \(C(G)\) .
\end{enumerate}

COROLLARY 2. Let C be a chamber of t and \(\overline{C}\) its closure. The canonical injections \(\overline{C} \rightarrow t \rightarrow g\) induce by passage to the quotient homeomorphisms

\[
\overline {{\mathrm{C}}} \rightarrow \mathfrak {t} / \mathrm{W} \rightarrow \mathfrak {g} / \operatorname{Ad} (\mathrm{G}).
\]

The canonical maps \(\overline{C} \rightarrow t\) and \(t \rightarrow t/W\) are proper (General Topology, Chap. III, §4, no. 1, Prop. 2 c)). The map \(\overline{C} \rightarrow t/W\) is continuous, proper and bijective (Prop. 2 c)); thus, it is a homeomorphism, hence the corollary in view of the Cor. of Prop. 6 of §2, no. 5.

Remark 4. Denote by \(\mathfrak{g}_{\mathrm{reg}}\) the set of regular elements of \(\mathfrak{g}\) (Chap. VII, §2, no. 2, Def. 2) and put \(\mathfrak{t}_{\mathrm{reg}} = \mathfrak{t} \cap \mathfrak{g}_{\mathrm{reg}}\) . For \(x \in \mathfrak{t}\) , we have

\[
\det (\mathrm{X} - \operatorname{ad} _ {\mathfrak {g}} x) = \mathrm{X} ^ {\dim t} \prod_ {\alpha \in \mathrm{R(G,T)}} (\mathrm{X} - \delta (\alpha) (x)),
\]

and hence \(t_{reg}\) is the set of elements x of t such that \(\delta(\alpha)(x) \neq 0\) for all \(\alpha \in \mathrm{R}(\mathrm{G}, \mathrm{T})\) , that is the union of the chambers of t (so \(t_r \subset t_{reg}\) ). Consequently \(\overline{C} \cap t_{reg} = C\) , so we have homeomorphisms

\[
\mathrm{C} \rightarrow \mathfrak {t} _ {\text { reg }} / \mathrm{W} \rightarrow \mathfrak {g} _ {\text { reg }} / \mathrm{Ad} (\mathrm{G}).
\]

COROLLARY 3. Assume that G is simply-connected; let g be a regular element of G. There exist a maximal torus S of G, and an alcove A of L(S), both uniquely determined, such that \(g \in \exp(A)\) and \(0 \in \overline{A}\) .

We can assume that \(g\) belongs to \(\mathrm{T}_r\) (§2, no. 2, Th. 2). Let \(x\) be an element of \(\mathfrak{t}_r\) such that \(\exp x = g\) , and let \(\mathrm{A}'\) be the alcove of \(\mathfrak{t}\) containing \(x\) . The alcoves \(\mathrm{A}\) of \(\mathfrak{t}\) such that \(g \in \exp(\mathrm{A})\) are the alcoves \(\mathrm{A}' - \gamma\) for \(\gamma \in \Gamma(\mathrm{T})\) ; thus, the assertion follows from Prop. 2 d).

